% ══════════════════════════════════════════════════════════════════════════════
\section{Theoretical Framework}
\label{sec:theory}
% ══════════════════════════════════════════════════════════════════════════════

Understanding the relationship between agricultural productivity and conflict requires engaging with several competing theoretical channels that generate opposing empirical predictions. This section outlines the three primary mechanisms---opportunity cost, predation, and resource scarcity---before discussing the institutional conditions that determine which mechanism dominates in a given context.

\subsection{The Opportunity Cost Channel}

The most widely cited theoretical link between income shocks and conflict runs through the opportunity cost of violence. Building on the economic model of rational crime \citep{Becker1968}, this framework treats individuals as agents who allocate effort between productive and destructive activities based on expected relative returns. When agricultural yields are high, the returns to farming and related labour-market activities increase, raising the cost of participation in conflict. Conversely, when harvests fail and rural income falls, the opportunity cost of violence declines, making recruitment into rebel movements, banditry, or communal raiding comparatively more attractive.

Formal models of conflict and appropriation \citep{Grossman1991} formalise this logic by modelling the allocation between production and predation as a function of the relative productivity of each activity. An adverse yield shock shifts the equilibrium allocation toward conflict by lowering the shadow price of violence relative to production. \citet{DalBoDalBo2011} extend this framework to show that both labour and capital income shocks affect conflict incentives, but through different mechanisms: a rise in labour productivity (as from a good harvest) increases the opportunity cost of conflict and reduces violence, while a rise in the value of appropriable capital goods (as from rising commodity prices) increases the prize from predation and thus conflict. This distinction is crucial for interpreting the sign of the yield--conflict relationship: the same positive income shock may reduce violence through opportunity cost effects while simultaneously increasing it through appropriation incentives.

\subsection{The Predation and Resource-Contestation Channel}

A second strand of theory predicts a positive relationship between agricultural productivity and conflict, operating through the incentive to appropriate surplus. Contest models of appropriation \citep{Tullock1980, Skaperdas1996, Hirshleifer1991} demonstrate that when the value of a contestable resource increases, the equilibrium investment in coercion also increases, as actors compete for the larger prize. Applied to agrarian settings, these models imply that abundant harvests expand the surplus available for extraction under conditions of weak property rights or limited enforcement capacity, thereby raising conflict risk.

Empirical evidence for the predation channel is provided by \citet{DubeVargas2013}, who show that positive price shocks to labour-intensive crops (coffee) reduce conflict in Colombia while positive price shocks to capital-intensive, easily appropriable commodities (oil) increase it---directly testing the distinction drawn by \citet{DalBoDalBo2011}. In agrarian Sub-Saharan Africa, where agricultural surpluses may be appropriated by armed groups through taxation, looting, or territorial control of productive land, a positive yield shock could generate more contestable resources and thus higher conflict, particularly in regions with weak state capacity.

\subsection{The Resource-Scarcity Channel}

A third mechanism emphasises competition under conditions of scarcity. When yields decline and agricultural resources become scarce, competition over land, water, and pastoral routes intensifies, particularly among groups whose livelihoods are most sensitive to marginal changes in resource availability. The environmental security literature \citep{HomerDixon1994} argues that resource scarcity interacts with population pressure and institutional weakness to generate instability, through pathways that include forced migration, breakdown of customary resource-sharing agreements, and heightened ethnic or communal tensions. Micro-level models of common-pool resource management \citep{Ostrom1990} further suggest that cooperative arrangements governing shared resources---such as communal grazing lands and water sources---become harder to sustain when scarcity increases.

This channel differs from the opportunity-cost channel in its emphasis on inter-group competition rather than individual-level participation decisions. Whereas the opportunity-cost framework predicts that negative yield shocks increase conflict by making individual violence more attractive relative to farming, the scarcity channel predicts that collective resource contests escalate as shared resources become more constrained.

\subsection{Institutional Mediation}

These three channels do not operate in isolation. The relative strength of each mechanism depends critically on the institutional environment, particularly the state's capacity to enforce property rights, adjudicate disputes, and provide alternatives to self-help in times of economic stress. \citet{BesleyPersson2009, BesleyPersson2011} develop a political-economy model in which the same economic shock can generate stability or conflict depending on whether fiscal and legal institutions are sufficiently strong to channel surplus into investment rather than violence and to buffer communities against adverse shocks through redistribution. In contexts with effective institutions, an abundant harvest generates investment and employment, raising opportunity costs and reducing conflict. In fragile states with weak enforcement, the same harvest creates appropriable surplus that invites predation.

The theoretical implication is therefore one of fundamental ambiguity: the net sign of the yield--conflict relationship is an empirical question, and heterogeneity in estimated effects across studies and contexts likely reflects genuine variation in the relative importance of each channel, conditioned by institutional quality and crop appropriability. This motivates the empirical approach taken in this paper, which seeks to estimate the reduced-form relationship between predicted yield and conflict across a diverse set of countries with varying degrees of state capacity and agricultural production structures.
